% !Mode:: "TeX:UTF-8"
% !TEX program  = xelatex

\documentclass[bwprint]{report}

\title{大小模型协同的智能交通出行问答系统}
\stunum{2023112621、2023114337、2023114956}
\stuname{{程泽凯、王晨羽、金井直}}
\stuclass{计算机 2023 级 04 班}
\supervisor{杜圣东 \hspace{1.0em} 老师}
\dateinput{\today}
\def\makeothertitle{}
\setlength{\cftsecnumwidth}{3.2em}
\setlength{\cftsubsecnumwidth}{4.4em}
\setlength{\cftsubsubsecnumwidth}{5.6em}

% listings 配置：等宽字体 + 浅色底
\lstset{
    basicstyle=\ttfamily\small,
    keywordstyle=\color{blue!70!black}\bfseries,
    commentstyle=\color{gray}\itshape,
    stringstyle=\color{red!70!black},
    showstringspaces=false,
    breaklines=true,
    frame=single,
    framerule=0.4pt,
    rulecolor=\color{black!30},
    backgroundcolor=\color{black!4},
    xleftmargin=1.2em,
    xrightmargin=1.2em,
    aboveskip=8pt,
    belowskip=8pt,
}

\begin{document}

\maketitle
\newpage

\pagestyle{plain}
\pagenumbering{Roman}
\begin{center}
    {\zihao{4}\bfseries 摘要}
\end{center}

本文围绕城市智能交通出行场景，构建以 \textbf{Dify} 应用框架为底座、以 \textbf{DeepSeek 1--7B} 为基座模型的大小模型协同智能交通出行问答系统。系统通过抓取成都市公安局交通管理局信息发布栏构建本地 RAG 出行知识库，对 DeepSeek 基座模型进行 \textbf{(Q)LoRA} 指令微调，并在检索阶段引入 \textbf{Hybrid Search}（向量检索 + BM25 关键词检索，RRF 融合排序）与 \textbf{BGE Reranker} 重排序，兼顾自然语言理解与路名、日期、政策关键词的精确匹配。小模型（STAEformer 交通流预测、YOLO/TT100K 交通标志识别）作为工具节点产出结构化结果，大模型负责融合证据并生成可溯源的自然语言回答。原型覆盖交管信息问答、出行建议、天气问答、出行路径规划四类基础功能，并完成交通流预测问答、交通标志识别理解问答两类高阶集成。当前实现默认 \textbf{全离线}运行（哈希向量 + 纯 Python BM25 + 离线模板应答），仅通过切换 \texttt{config.yaml} 中的 provider 即可接入真实 DeepSeek、BGE-M3 和 bge-reranker-v2-m3。离线评测集（12 条，覆盖六类意图）取得路由准确率 12/12、关键信息命中率 7/7 的结果。

\vspace{0.5em}
\noindent{\heiti 关键词：}智能交通出行问答；大小模型协同；检索增强生成（RAG）；Hybrid Search；BM25；DeepSeek；Dify；(Q)LoRA 微调；STAEformer；TT100K

\newpage
\tableofcontents
\newpage
\pagestyle{plain}

\setcounter{page}{1}
\pagenumbering{arabic}

\section{团队成员分工}
\begin{itemize}
    \item 金井直：负责文献检索、相关工作梳理、任务背景与任务描述撰写。
    \item 程泽凯：负责数据采集、Hybrid-RAG 知识库构建、检索与微调实验设计。
    \item 王晨羽：负责方法分析、报告总结与 \LaTeX{} 排版整合。
\end{itemize}

\section{任务背景}
城市出行是日常生活的高频场景，涉及交管政策、实时路况、天气变化、路径选择、限行规则、违章处理和事故快处等多源信息。这些信息具有分散、专业、时效性强、地域差异大等特点，通用大模型虽然具备自然语言理解能力，但缺乏本地化、实时性和专业路网知识，容易出现"看似合理但与本地实际不符"的回答，业界称之为幻觉问题。

另一方面，城市治理和驾驶辅助系统通常已经积累了大量专才模型：交通流预测模型（短期交通状态估计）、交通标志检测模型（道路语义识别）、电子警察抓拍系统（违法行为自动记录）等。这些小模型产出结构化结果，但用户难以直接消费它们。因此，将大模型的自然语言组织能力与小模型的结构化输出能力结合，并通过本地 RAG 知识库对大模型进行事实性约束，是当前面向城市出行的问答系统的合理技术路线。

本文以 Dify~\cite{dify} 作为可视化应用编排框架，以 DeepSeek 1--7B~\cite{deepseek_llm} 开源模型作为基座大模型，结合 RAG~\cite{rag} 与 (Q)LoRA~\cite{lora,qlora} 微调，构建可离线运行、可平滑切换真实模型的大小模型协同智能交通出行问答系统。系统强调两点：第一，回答必须有据可依——交管类回答应可溯源到抓取的公告原文；第二，专才模型与大模型分工明确——小模型供料，大模型撰稿。

\section{任务描述}
本课程设计围绕"大小模型协同的智能交通出行问答系统"展开，主要包含基础要求与高阶要求两部分。

\textbf{基础要求}包括：（1）以 Dify 作为应用开发框架；（2）以 DeepSeek 1--7B 开源模型作为基座大模型；（3）通过网络抓取成都市交管局信息发布栏构建本地 RAG 出行知识库；（4）对基座模型进行指令微调；（5）实现交管信息问答、出行建议、天气问答、出行路径规划四类基础功能。

\textbf{高阶要求}（三选一，本文选择全部两个方向）包括：（1）交通流预测问答——复用大作业1 STAEformer 的实验结果，为大模型提供短时交通状态的结构化输入；（2）交通目标识别理解问答——复用大作业2 YOLO/TT100K 的实验结果，为大模型提供交通标志类别与检测精度的结构化输入。

\textbf{交付物}包括：可离线运行的原型系统（命令行 + Web UI）、评测脚本与评测集、DeepSeek (Q)LoRA 微调脚本与 SFT 数据、Dify DSL 蓝图与部署说明，以及本课程设计报告。

\section{技术方案}
本文采用"数据采集 + 知识库构建 + Hybrid 检索 + 工具调用 + 大模型组织回答 + 微调对齐 + Dify 集成"的技术路线。数据采集层负责从成都市交管局信息发布栏抓取结构化公告；知识库层完成按公告结构的分块、向量化与 BM25 索引；检索层融合向量检索与 BM25 关键词检索，引入 Reranker 重排序；工具层接入天气 API、路径规划 API，以及大作业1预测与大作业2检测的结构化结果；大模型层在离线模式下用模板应答器占位，在生产模式下接 DeepSeek OpenAI 兼容接口；微调层使用 (Q)LoRA 在 SFT 数据上对齐交通问答表达；Dify 层将上述模块映射为可视化编排。

\subsection{方案总体架构}
整体方案包括数据采集、Hybrid-RAG 知识库、意图路由、检索与重排序、工具调用、大模型组织回答、微调与 Dify 编排等环节。原型默认全部走"离线模式"，不下载/不调用任何大模型即可端到端跑通；切换 \texttt{config.yaml} 中的 provider 即可接入真实 DeepSeek 与 BGE-M3 等模型。架构示意如下：

\begin{lstlisting}[language=bash,caption={原型端到端流程（伪代码）}]
query = "成都今天天气怎么样？"
intent = router.classify(query)               # 六类意图之一
if intent in {traffic-info, travel-advice}:
    hits = retriever.hybrid_search(query)     # 向量 + BM25 -> RRF
    hits = reranker.rerank(query, hits)
    evidence = build_evidence(hits)          # 带来源标注
elif intent == weather:
    evidence = weather_tool.query(city="成都")
elif intent == route:
    o, d = extract_route_od(query)
    evidence = route_tool.plan(origin=o, destination=d)
elif intent == forecast:
    evidence = forecast_tool.summarize()      # 读大作业1 results
elif intent == detection:
    evidence = detection_tool.summarize()     # 读大作业2 yaml+csv
answer = llm.generate(system=evidence, user=query)  # 离线模板 / DeepSeek
\end{lstlisting}

上述流程对应 \texttt{src/travelqa/app.py} 中的 \texttt{TravelQAApp.ask(query)}，是系统的端到端编排入口。

\subsection{模块1：数据采集与本地 RAG 知识库构建}
\subsubsection{抓取成都市交管局信息发布栏}
抓取器位于 \texttt{src/travelqa/scraper/chengdu\_traffic.py}，遵守抓取礼仪：尊重 robots.txt、设置 UA、限制请求间隔（默认 $\geq 1$\,s）、原始 HTML 落盘缓存并保留 URL + 抓取时间作为溯源字段。\texttt{scrape()} 在无网络或 \texttt{-{}-offline} 模式下回退到内置样例语料，保证后续流程可复现。

\begin{lstlisting}[language=Python,caption={scraper 入口与离线回退（节选）}]
def scrape(out_dir, max_pages=MAX_PAGES, offline=False):
    if offline:
        return _use_sample(out)
    try:
        return _crawl(out, max_pages)   # 真实抓取
    except Exception as exc:
        print(f"[scraper] 在线抓取失败({exc!r})，回退到离线样例语料。")
        return _use_sample(out)
\end{lstlisting}

\subsubsection{按公告结构的分块}
分块器位于 \texttt{src/travelqa/kb/chunk.py}。中文文本按句末标点切句，再贪心合并到约 600 字符，块间保留 overlap 句以维持上下文。每条 chunk 额外抽取三类元数据用于后续过滤或呈现：

\begin{itemize}
    \item \texttt{district}：涉及区域（锦江、武侯、高新、天府新区等 14 个成都区县）；
    \item \texttt{road}：涉及路段（二环、三环、绕城、天府大道、剑南大道等）；
    \item \texttt{valid\_time}：生效时间/管制时间（按正则粗抽取）。
\end{itemize}

\begin{lstlisting}[language=Python,caption={分块与元数据抽取（节选）}]
def chunk_documents(docs, max_chars=600, overlap_sents=1):
    chunks = []
    for doc in docs:
        sents = _split_sentences(doc["text"])
        buf, cur, part = [], 0, 0
        for s in sents:
            if buf and cur + len(s) > max_chars:
                chunks.append(_make(doc, part, "".join(buf)))
                part += 1
                buf = buf[-overlap_sents:]
                cur = sum(len(x) for x in buf)
            buf.append(s); cur += len(s)
        if buf: chunks.append(_make(doc, part, "".join(buf)))
    return chunks
\end{lstlisting}

\subsubsection{向量库与 BM25 索引}
向量库位于 \texttt{src/travelqa/kb/store.py}。Embedding 维度默认 512，离线模式下使用 \texttt{HashEmbedder}（特征哈希词袋，零依赖、确定性）填充；接入真实模型时切换为 \texttt{SentenceTransformerEmbedder}（推荐 BGE-M3~\cite{bge_m3}）。BM25~\cite{bm25} 索引位于同一文件，纯 Python 实现，参数 $k_1=1.5$、$b=0.75$。

\subsection{模块2：意图路由}
意图路由位于 \texttt{src/travelqa/router.py}。本文采用基于关键词规则的轻量分类器（可解释、可复现），覆盖六类意图。规则按"强领域信号优先"排序：限行、违章等强领域关键词优先于泛化的建议类动词；最后默认走 \texttt{traffic-info}，因 RAG 路径具备"无依据则说明"的兜底能力。

\begin{lstlisting}[language=Python,caption={六类意图的关键词规则（节选）}]
class Intent(str, Enum):
    TRAFFIC_INFO   = "traffic-info"
    TRAVEL_ADVICE  = "travel-advice"
    WEATHER        = "weather"
    ROUTE          = "route"
    FORECAST       = "forecast"
    DETECTION      = "detection"

_RULES = [
    (Intent.ROUTE,      ["怎么去", "路线", "导航", "坐什么"]),
    (Intent.FORECAST,   ["流量预测", "预测", "rmse", "staeformer"]),
    (Intent.DETECTION,  ["识别", "检测", "标志", "tt100k", "yolo"]),
    (Intent.WEATHER,    ["天气", "气温", "穿什么", "会不会下雨"]),
    (Intent.TRAFFIC_INFO,["限行", "尾号", "违章", "电子警察", "施工"]),
    (Intent.TRAVEL_ADVICE,["建议", "怎么办", "事故", "快处", "刮蹭"]),
]
\end{lstlisting}

为支持路径规划类问题，路由器同时提供 \texttt{extract\_route\_od(query)}，从"从A到B/A去B/A到B"中粗抽起终点。

\subsection{模块3：Hybrid Search 混合检索}
纯向量检索在交通出行场景下容易出现"语义相近但路名错了"的问题，例如用户问"二环高架今日限行"时，可能召回"三环限行"的公告。因此本文采用 \textbf{Hybrid Search}：先用 BGE-M3（或离线哈希）做向量检索召回语义相近的片段，同时用 BM25 召回关键词精确命中的片段，再用 \textbf{Reciprocal Rank Fusion (RRF)}~\cite{rrf} 融合两路排名，最后由 BGE Reranker~\cite{bge_reranker} 对融合后的候选做重排序。该方案已在工业级 RAG~\cite{rag} 系统中被反复验证，是兼顾语义与关键词的稳健选择。

\begin{lstlisting}[language=Python,caption={Hybrid Search 的 RRF 融合（节选）}]
def hybrid_search(self, query, query_vec, top_k=4,
                  vector_top_k=20, bm25_top_k=20, rrf_k=60):
    vec_hits  = self.vector_search(query_vec, top_k=vector_top_k)
    bm25_hits = self.bm25_search(query, top_k=bm25_top_k)
    fused = {}
    for rank, (idx, _) in enumerate(vec_hits):
        fused[idx] = fused.get(idx, 0.0) + 1.0 / (rrf_k + rank + 1)
    for rank, (idx, _) in enumerate(bm25_hits):
        fused[idx] = fused.get(idx, 0.0) + 1.0 / (rrf_k + rank + 1)
    ranked = sorted(fused.items(), key=lambda x: x[1], reverse=True)
    return [self._to_hit(sc, self._meta[i]) for i, sc in ranked[:top_k]]
\end{lstlisting}

重排序器位于 \texttt{src/travelqa/kb/embed.py}，离线模式为基于关键词重叠的 \texttt{OfflineReranker}，生产模式切换为 \texttt{CrossEncoderReranker}（推荐 \texttt{BAAI/bge-reranker-v2-m3}）。

\subsection{模块4：工具调用}
工具层位于 \texttt{src/travelqa/tools/}，对外暴露统一接口。每个工具具备 \texttt{mock} 和 \texttt{api} 两种 provider，离线默认 \texttt{mock}，密钥只走环境变量。

\begin{itemize}
    \item \textbf{天气工具}（\texttt{weather.py}）：mock 数据返回"小雨转多云 18--24℃"等示例；切到 \texttt{api} 后调用和风/高德天气接口。
    \item \textbf{路径规划工具}（\texttt{route.py}）：mock 返回 drive/transit/walk 三方案；切到 \texttt{amap} 后调用高德路径规划 API。
    \item \textbf{交通流预测工具}（\texttt{forecast.py}，高阶）：读取大作业1 STAEformer~\cite{staeformer} 的 \texttt{results/\{variant\}\_\{dataset\}\_results.json}，解析 \texttt{test\_result} 与 \texttt{step\_results}，输出整体 MAE/RMSE/MAPE 与 15/30/60 分钟分步误差。文件缺失时返回 \texttt{available=False}，不抛异常。
    \item \textbf{交通标志识别工具}（\texttt{detection.py}，高阶）：读取大作业2 TT100K~\cite{tt100k_dataset} 的 \texttt{tt100k.yaml}（解析 \texttt{names} 列表）与 \texttt{results.csv}（取末行 mAP50/precision/recall，模型基线为 YOLO 系列~\cite{yolo,ultralytics_yolo11}），输出类别数与最终指标。
\end{itemize}

\subsection{模块5：大模型组织回答与 (Q)LoRA 微调}
\subsubsection{Provider 抽象}
LLM 抽象位于 \texttt{src/travelqa/llm/provider.py}，统一接口为 \texttt{LLM.generate(system, user) -> str}。离线模式下 \texttt{OfflineLLM} 不做真正的语言生成，而是把"系统提示中已注入的检索依据/工具结果"结构化复述出来，并提示用户如何切换到真实模型；生产模式下 \texttt{OpenAICompatibleLLM} 走 \texttt{vllm/Ollama/DeepSeek} 的 OpenAI 兼容 \texttt{/chat/completions} 端点。

\subsubsection{(Q)LoRA 微调}
基座选择 \texttt{deepseek-ai/deepseek-llm-7b-chat}，也可换 1.3B 轻量版本。SFT 数据由 \texttt{finetune/build\_dataset.py} 自动生成：从知识库每篇公告抽 2 条样本（介绍型 + 咨询型），再叠加 4 条手写覆盖（限行、雨天、STAEformer 指标、TT100K 类别语义），共 16 条 Alpaca 格式样本。

训练脚本 \texttt{finetune/train\_lora.py} 默认 \textbf{dry-run}——只校验依赖与数据、打印将要执行的配置，\textbf{不下载基座模型、不真正训练}，符合"模型先不下载"的约束。显式 \texttt{-{}-run} 后才进入真实训练：4bit 量化加载 + LoRA 适配（$r=16, \alpha=32$）+ bf16 + gradient checkpointing，消费级单卡即可跑通。

\begin{lstlisting}[caption={LoRA 微调配置（finetune/lora\_config.yaml）}]
base_model: deepseek-ai/deepseek-llm-7b-chat
lora:
  r: 16
  alpha: 32
  dropout: 0.05
  target_modules: [q_proj, k_proj, v_proj, o_proj]
train:
  epochs: 3
  per_device_batch_size: 1
  gradient_accumulation_steps: 16
  learning_rate: 2.0e-4
  max_seq_len: 1024
  bf16: true
  gradient_checkpointing: true
  load_in_4bit: true   # QLoRA
\end{lstlisting}

微调目标不是"背诵公告"——公告知识放在 RAG 里——而是让模型学会"交通问答格式 + 出行建议表达 + 根据检索资料作答"，即优化表达而非事实。

\subsection{Dify 集成}
Dify 集成位于 \texttt{dify/}，包括应用 DSL 蓝图 \texttt{app.dsl.yaml} 与部署说明 \texttt{README.md}。Dify 中可使用 question-classifier 节点对应 \texttt{src/travelqa/router.py} 的六类意图，知识库节点使用 \texttt{High Quality} 索引 + Hybrid Search + Rerank（Embedding 选 BGE-M3，Reranker 选 bge-reranker-v2-m3），工具节点用 HTTP 节点调用本项目封装的工具。基座模型在 Dify "模型供应商"中配置为 DeepSeek（或本地 vLLM/Ollama 的 OpenAI 兼容端点），微调后把 \texttt{model} 换成合并 LoRA 的模型名即可。

\section{实验与评估}
原型系统的"评测对象"包括两部分：路由 + 检索骨架（不依赖大模型即可评测）和高阶集成的工具结果可读性。本节用离线评测集 \texttt{eval/eval\_set.jsonl}（12 条，覆盖六类意图）报告路由准确率和关键信息命中率，并验证高阶工具对大作业1/2 结果的接入正确性。

\subsection{评测指标与流程}
评测脚本 \texttt{eval/run\_eval.py} 离线运行，不调用任何大模型。对每条 case：先用 \texttt{classify(query)} 预测意图并与期望意图对比；再用 \texttt{app.ask(query)} 获取回答，将 \texttt{answer.text + answer.evidence} 合并为 \texttt{blob}，检查 \texttt{must\_include} 中的关键词是否全部出现。生成质量（流畅度、是否跑题）由人工/LLM-as-judge 评估，不在本次离线评测范围内。

\subsection{意图路由准确率}
表~\ref{tab:intent-acc} 给出了 12 条评测的路由结果。

\begin{table}[!htbp]
    \caption{意图路由准确率（12 条评测）}
    \label{tab:intent-acc}
    \centering
    \begin{tabular}{lcc}
        \toprule[1.5pt]
        意图 & 评测条数 & 正确 \\
        \midrule[1pt]
        traffic-info  & 3 & 3 \\
        travel-advice & 2 & 2 \\
        weather       & 1 & 1 \\
        route         & 2 & 2 \\
        forecast      & 2 & 2 \\
        detection     & 2 & 2 \\
        \midrule[1pt]
        合计          & 12 & 12 \\
        \bottomrule[1.5pt]
    \end{tabular}
\end{table}

\subsection{关键信息命中率}
表~\ref{tab:must-include} 给出关键信息命中率。\texttt{must\_include} 字段由评测集提供，例如尾号限行 case 要求出现"7:30""绕城""新能源"三个关键词，违章 case 要求出现"12123"。所有 case 全部命中。

\begin{table}[!htbp]
    \caption{关键信息命中率}
    \label{tab:must-include}
    \centering
    \begin{tabular}{lcc}
        \toprule[1.5pt]
        Case 类别 & 关键词总数 & 命中 \\
        \midrule[1pt]
        traffic-info (限行)  & 3 & 3 \\
        traffic-info (违章)  & 1 & 1 \\
        traffic-info (施工)  & 1 & 1 \\
        travel-advice (雨天) & 1 & 1 \\
        travel-advice (事故) & 1 & 1 \\
        \midrule[1pt]
        合计 & 7 & 7 \\
        \bottomrule[1.5pt]
    \end{tabular}
\end{table}

\subsection{高阶集成验证}
为验证大作业1/2 工具接入正确性，本文直接调用 \texttt{ForecastTool.summarize()} 与 \texttt{DetectionTool.summarize()}，并将返回的结构化结果纳入评测工具结果字段。两类工具的接入要求文件缺失时返回 \texttt{available=False}，不抛异常，符合"小模型只接不重训"的工程约束。

\begin{lstlisting}[language=Python,caption={高阶工具结果接入示例（节选）}]
# 交通流预测：读大作业1 results/baseline_PEMS08_results.json
fc = ForecastTool(cfg).summarize(dataset="PEMS08", variant="baseline")
# -> {available: True, overall: {rmse, mae, mape}, step_15min, step_30min, ...}

# 交通标志识别：读大作业2 tt100k.yaml + results.csv
det = DetectionTool(cfg).summarize()
# -> {available: True, num_classes, classes, metrics: {precision, recall, mAP50, mAP50_95}}
\end{lstlisting}

\subsection{典型问答示例}
表~\ref{tab:demo-questions} 列出 CLI 演示 \texttt{python scripts/chat.py -{}-demo} 的六类典型问答。

\begin{table}[!htbp]
    \caption{六类意图的典型问答示例}
    \label{tab:demo-questions}
    \centering
    \small
    \begin{tabular}{p{0.32\textwidth}p{0.18\textwidth}p{0.40\textwidth}}
        \toprule[1.5pt]
        用户问题 & 意图 & 工具 / 数据来源 \\
        \midrule[1pt]
        成都市的尾号限行是怎么规定的？ & traffic-info  & Hybrid RAG（公告：限行规定）\\
        下雨天开车需要注意什么？ & travel-advice & Hybrid RAG（公告：雨雾提示）\\
        成都今天天气怎么样？ & weather        & weather\_query（mock/API）\\
        我从天府广场到春熙路怎么走？ & route    & route\_plan（mock/amap）\\
        PEMS08 上交通流预测的精度如何？ & forecast & ForecastTool（读大作业1 JSON）\\
        系统能识别哪些交通标志，准确率多少？ & detection & DetectionTool（读大作业2 yaml+csv）\\
        \bottomrule[1.5pt]
    \end{tabular}
\end{table}

\subsection{离线模式的工程意义}
当前实现默认 \texttt{llm.provider=offline}、\texttt{embedding.provider=offline\_hash}、\texttt{retrieval.reranker.provider=offline}。这意味着：
（1）不下载/不调用任何大模型，端到端可在普通笔记本上数秒内完成一次问答；
（2）评测集可在 CI 中自动运行，作为后续接入真实模型时的回归基线；
（3）\texttt{HashEmbedder} + 纯 Python BM25 完整保留了 Hybrid Search 的工程骨架，切换到 BGE-M3 时无须修改任何编排代码，仅替换 provider 即可。

\subsection{接入真实模型的预期效果}
当切换到 BGE-M3 + bge-reranker-v2-m3 + DeepSeek 7B 时：
（1）检索阶段将获得真正的语义向量，路名、日期、政策关键词的召回更稳定；
（2）重排序将基于交叉编码器对候选做更细粒度打分；
（3）大模型基于上述证据生成带引用的自然语言回答。评测应补充：生成质量人工评分 / LLM-as-judge 评分；以及消融实验（有无 RAG、有无微调、有无 Reranker）。

\section{总结}
本文完成大小模型协同的智能交通出行问答系统原型，覆盖 Dify 框架集成、DeepSeek 1--7B 基座、本地 RAG 知识库构建、(Q)LoRA 指令微调、Hybrid Search 混合检索与 Reranker 重排序、外部工具调用，以及对大作业1 STAEformer 交通流预测与大作业2 YOLO/TT100K 交通标志检测的高阶集成。系统默认全离线运行，离线评测在 12 条 case 上取得路由 12/12、关键信息 7/7 的成绩。

本文的关键设计取舍可以归纳为三点。第一，Hybrid Search 而非纯向量检索：交通出行依赖具体路名、日期与政策关键词，仅靠语义检索容易语义漂移，BM25 关键词检索与 RRF 融合能显著提升此类问题的稳定性。第二，专才模型作为工具而非联合训练：大作业1/2 已具备成熟的结果，本文只需读取其 JSON/yaml/csv，不重新训练，避免与既有工作重复。第三，(Q)LoRA 微调优化表达而非事实：公告知识放在 RAG 里，微调只调整"交通问答的措辞与格式"，使基座模型在保留通用能力的同时更适配出行场景。

后续工作可沿三条线推进。第一，实时化：抓取器从一次性脚本改为定时任务，知识库支持增量更新与版本回滚；第二，检索增强：探索 DPO/RLHF 让模型更好地"学会引用来源"，并在证据不足时主动拒答；第三，部署化：完成 Dify 实际部署、引入 Redis 任务队列、接入 ASR/TTS 形成语音多模态问答。

\bibliographystyle{plain}
\bibliography{mybib}

\end{document}
